\clearpage
\SourceLeaf{276}
\renewcommand{\BookRunningTitle}{The Periods of Classes}
\pdfbookmark[0]{Proof of Some Theorems Concerning the Periods of Classes}{periodes-classes}
\begin{center}
\Large\bfseries PROOF OF SOME THEOREMS CONCERNING\par\medskip
\large THE PERIODS OF CLASSES OF BINARY QUADRATIC FORMS.
\end{center}
\BookRule{.18\textwidth}

\pdfbookmark[1]{Theorem I}{periodes-theoreme-I}
\textsc{Theorem. I.}\quad\emph{The number of (properly primitive) classes of the same determinant which, when raised to the $P^{\mathrm{th}}$ power, where $P$ is either a prime number or a prime power $=p^\pi$, produce the principal class $K$, is equal either to $1$ or to a power of this same prime number $p$.}

\emph{Proof.}\quad Let $(\mathfrak Q)$ be the whole group of all the classes in question, and $n$ their number. Since the principal class $K$ is necessarily contained in $(\mathfrak Q)$, the theorem is evident if it is the only one there. But if there are others, the number of classes contained in the period of each will be a power of $p$; let one of them be $A$, and suppose that its period $(\mathfrak A)$ contains $p^\alpha$ classes, all of which will be included in $(\mathfrak Q)$. Now, if the classes of this period $(\mathfrak A)$ exhaust $(\mathfrak Q)$, we shall have $p^\alpha=n$, and the theorem will be proved; otherwise, let $B$ be any class of $(\mathfrak Q)$ not contained in $(\mathfrak A)$, and suppose its period is developed until we reach a class $bB$ which is also among the classes of $(\mathfrak A)$. This must necessarily happen, because at least the principal class is common to this period and to $(\mathfrak A)$. Now, supposing that $bB$ is the \emph{first} class in the period of $B$ common to $(\mathfrak A)$, or that $b$ is as small as possible, I claim

$1^0$.\quad That $b$ will be a power of $p$. For it is evident that, on putting $b=p^\beta h$, $bB=iA$ and $hk\equiv1\pmod{p^\pi}$ (which will be possible), we shall have $kbB=p^\beta hkB=p^\beta B=ikA$,
\SourceLeaf{277}
that is, $p^\beta B$ will also be among the classes of $(\mathfrak A)$, from which it follows that $h=1$ and $b=p^\beta$.

$2^0$.\quad That, if the classes $K,B,2B\dotsfour(b-1)B$ are denoted by $(\mathfrak B)$, all the compositions of a class of $(\mathfrak A)$ with a class of $(\mathfrak B)$ will give $p^{\alpha+\beta}$ different classes. For, supposing $mA+nB=m'A+n'B$ and $n=n'$, we shall necessarily have $m=m'$; if $n>n'$, we shall have $(n-n')B=(m'-m)A$, which is impossible unless $n=n'$.

$3^0$.\quad That these $p^{\alpha+\beta}$ different classes will be included in $(\mathfrak Q)$, which is evident.

Now, if these $p^{\alpha+\beta}$ classes exhaust $(\mathfrak Q)$, the theorem is proved; otherwise, we shall choose another class of $(\mathfrak Q)$ not contained among them, namely $C$; we shall continue its period until we reach a class already included among the classes composed from $(\mathfrak A)$ and $(\mathfrak B)$. By reasoning similar to the preceding\textsuperscript{'}, we shall prove that the exponent of this class must be a power of $p$, $=p^\gamma$, and that the composition of the first $p^\gamma$ classes of the period of $C$ with the $p^{\alpha+\beta}$ classes already found will give $p^{\alpha+\beta+\gamma}$ different classes, all included in $(\mathfrak Q)$. If these classes still do not exhaust $(\mathfrak Q)$, we shall treat a fourth class $D$, etc., in the same way; and it is evident that, since $(\mathfrak Q)$ consists of a finite number of classes, these operations will also come to an end, and $n$ will be equal to a power of $p$.\quad Q. E. D.

\par\bigskip
\pdfbookmark[1]{Theorem II}{periodes-theoreme-II}
\textsc{Theorem. II.}\quad\emph{If the number of all the classes of the principal genus is expressed by $a^\alpha b^\beta c^\gamma$, etc., where $a,b,c$ denote different prime numbers, there will be in this genus $a^\alpha,b^\beta,c^\gamma$, etc., classes which, when raised respectively to the powers $a^\alpha,b^\beta,c^\gamma$, etc., produce the principal class.}

\emph{Proof}\quad Let $A,A',A''$, etc., be all the classes which, when raised to the power $a^\alpha$, produce $K$, and let $(\mathfrak A)$ be their totality; likewise $B,B',B''$, etc., $(\mathfrak B)$, $C,C',C''$, $(\mathfrak C)$, etc., etc. I claim that the composition of all the classes of $(\mathfrak A)$ with all the classes of $(\mathfrak B)$ with all the classes of $(\mathfrak C)$, etc., will produce classes different from one another. For if $A+B+C\dots=A'+B'+C'\dots$, etc., then, on putting $A-A'=A''$, $B-B'=B''$, etc., we shall have
\[
A''+B''+C''\text{ etc.}=K
\]
and hence, on raising to the power $b^\beta c^\gamma$, etc., $(b^\beta c^\gamma\dots)A''=K$, from which it follows readily that
\SourceLeaf{278}
$A''=K$ and $A=A'$, and in the same way we shall have $B=B'$, $C=C'$, etc. Let the totality of these classes be $=(S)$. Further, it is clear that all these classes will belong to the principal genus. Finally, there can be no class in the principal genus which is not included in $(S)$. Let\ .\ .\ .
\BookRule{.34\textwidth}

\par\bigskip
\pdfbookmark[1]{Remark}{periodes-bemerkung}
\begin{center}\textsc{Remark.}\end{center}
\begingroup\small
This fragment, written in 1801, relates to Disq. Arithm., Article~306, \textsc{ix}. The word \emph{dignité} is used here in a sense that is not otherwise customary.
\begin{flushright}\textsc{Stern}.\end{flushright}
\endgroup
\BookRule{.34\textwidth}
